\documentclass[conference]{IEEEtran}
\IEEEoverridecommandlockouts
\usepackage[T1]{fontenc}
\usepackage{cite}
\usepackage{amsmath,amssymb,amsfonts}
\usepackage{booktabs}
\usepackage{array}
\usepackage{graphicx}
\usepackage{url}
\usepackage{xcolor}
\usepackage{balance}
\begin{document}

\title{A Benchmark for Closed-Loop Human-Machine Integration}

\author{\IEEEauthorblockN{Survey, Idea, ExperimentDesign, and Author Agents}
\IEEEauthorblockA{Research synthesis and protocol proposal\\
Human-Machine Systems Study\\
DESIGN\_ONLY artifact}}

\maketitle

\begin{abstract}
The prospect of integrating people with computers is often described as the creation of a human-machine hybrid species. That phrase combines technologies with very different mechanisms and risks. This paper converts it into a testable engineering question: can a human and a computer form a stable, bidirectional, closed-loop system that improves a bounded task while preserving agency, biological safety, privacy, cybersecurity, and safe disengagement? We survey neural prostheses, non-invasive physiological interfaces, wearable assistance, and adaptive control, then propose the CLOSED-LOOP-HUMAN-MACHINE-INTEGRATION-BENCHMARK. The benchmark compares interface level, sensory feedback, human and machine adaptation, and governance under matched task demands. Its protocol proceeds from simulation to hardware-in-the-loop and only then to ethics-reviewed low-risk human validation. Primary measures include task utility, perturbation recovery, veto success, decoder drift, machine-off transfer, emergency-stop latency, privacy leakage, and attack success. The paper is a design proposal, not a report of completed experiments. It does not claim that a new biological species, sentient machine, or altered human moral status has been demonstrated.
\end{abstract}

\begin{IEEEkeywords}
brain-machine interface, human-machine integration, adaptive control, agency, neurotechnology, safety, privacy
\end{IEEEkeywords}

\section{Introduction}

Smart prostheses, exoskeletons, implantable sensors, artificial intelligence, and genomic editing have made human-machine integration a plausible research topic. Yet the question ``Could we integrate with computers to form a human-machine hybrid species?'' is scientifically under-specified. A wearable robot coupled through force and motion is not equivalent to a neural decoder. A neural prosthesis that translates motor-cortex activity into a robotic-arm command is not equivalent to a memory aid or a cognitive co-pilot. Germline modification is a separate biological and governance problem rather than a necessary step in building a useful interface.

This paper therefore adopts a narrower and falsifiable target. We ask whether a person and a computer can form a stable functional coupling in which information flows from the person to the machine and, where appropriate, back to the person; both sides adapt; performance improves on held-out conditions; and the person retains agency, safety, privacy, and the ability to disengage. This reframing is consistent with established neural-prosthesis evidence. People with tetraplegia have controlled multidimensional reach and grasp of a robotic arm through a motor-cortex array \cite{hochberg2012}, and earlier work demonstrated neuronal ensemble control of prosthetic devices by a human participant \cite{hochberg2006}. These are important demonstrations of task-specific coupling. They are not evidence of biological speciation.

The contribution is a benchmark and staged experiment design rather than a claim of a completed hybrid. The benchmark separates six integration levels, proposes a controlled factorial matrix, formalizes coupled-system observables, and treats agency, safe shutdown, privacy, and cybersecurity as primary engineering properties. It also specifies a machine-off transfer test to distinguish durable adaptation from short-term assistance. The protocol follows a `DESIGN\_ONLY` boundary: no human, animal, implant, cognitive, or genomic results are invented.

The rest of the paper is organized as follows. Section II summarizes the survey and defines integration levels. Section III states the benchmark idea and hypotheses. Section IV gives the closed-loop model and metrics. Section V presents the staged experiment design. Section VI describes statistical analysis and reproducibility. Section VII addresses safety, governance, and ethical boundaries. Section VIII discusses limitations and expected interpretations, followed by the conclusion.

\section{Survey: What Counts as Integration?}

\subsection{A technology ladder}

The word integration hides a ladder of increasing coupling and risk. We use the following operational levels.

\begin{table}[t]
\caption{Operational levels of human-machine integration}
\label{tab:levels}
\centering
\small
\begin{tabular}{p{0.08\columnwidth}p{0.27\columnwidth}p{0.47\columnwidth}}
\toprule
Level & Interface & Testable function and principal concern\\
\midrule
L0 & Wearable physical assistance & Forces and motion are exchanged externally; test stability, effort, fatigue, and aftereffects.\\
L1 & Non-invasive physiological & Electromyography, electroencephalography, inertial or eye signals control a device; test noise, drift, and privacy.\\
L2 & Implanted unidirectional & Implanted neural activity is decoded into commands; test useful control, stability, clinical risk, and maintenance.\\
L3 & Bidirectional embodied & Commands and sensory feedback form a sensorimotor loop; test prediction, plasticity, agency, and safe override.\\
L4 & Cognitive augmentation & A system supports memory, attention, language, or decisions; test mental autonomy and unintended inference.\\
L5 & Genomic or heritable & Biological substrates are modified; this is a distinct and higher-risk problem, excluded from the initial interface study.\\
\bottomrule
\end{tabular}
\end{table}

L0 and L1 can be studied with reversible, low-risk systems. L2 has already achieved meaningful task demonstrations but depends on surgery, electrode-tissue interfaces, calibration, and long-term signal stability. L3 is attractive because feedback may let a user predict device state and incorporate it into control. L4 changes the target from movement to cognition and should not inherit the evidence standards of a prosthetic-arm study. L5 is not a computer-interface extension: it raises heritable effects, consent across generations, and population-level governance. Table \ref{tab:levels} prevents these distinct claims from being collapsed into one binary question.

\subsection{Evidence from neural prostheses}

Hochberg et al. reported reach and grasp by people with tetraplegia using a neurally controlled robotic arm \cite{hochberg2012}. The study used a 96-channel motor-cortex array and demonstrated functional multidimensional control, including an action requiring reach and grasp. The result supports three modest conclusions: motor-related neural activity can be mapped to useful commands; a decoder can remain useful for a period of years in at least one implanted context; and an external machine can participate in an action loop with a person. It does not establish a general-purpose merger of human and computer cognition.

Earlier neuronal-ensemble work showed that neural activity could control prosthetic devices \cite{hochberg2006}. Reviews have emphasized open challenges involving signal stability, calibration, channel count, feedback, and practical deployment \cite{gilja2011}. A foundational review also describes brain-machine interfaces as a field spanning neural recording, decoding, learning, and actuation \cite{lebedev2006}. Together, these results motivate a benchmark that evaluates an entire coupled loop, rather than only the classifier accuracy of its decoder.

\subsection{Wearable and non-invasive routes}

Wearable robots produce a different kind of coupling. A controller senses gait phase, posture, or intent and applies bounded force while the person changes their motor policy. The central outcome is not maximum actuator torque. It is whether assistance remains useful across fatigue, terrain, delay, and perturbation without destabilizing the user or preventing independent operation. A non-invasive physiological interface has lower surgical risk but often lower signal-to-noise ratio, greater inter-session drift, and large individual variability. Comparing these routes with implanted systems requires common tasks and risk-adjusted outcomes, not a throughput contest.

\subsection{The missing benchmark}

The survey identifies seven recurring gaps. Acute control accuracy is more common than long-term co-adaptation. Wearable assistance, neural control, and cognitive augmentation are discussed as if they had the same observables. Agency and successful veto are under-measured. Studies rarely ask whether a user can operate after the machine is switched off. Neural-data privacy is often separated from control evaluation, and cybersecurity is not consistently tested as a biological-safety property. Finally, genomic editing is sometimes presented as part of hybridization even though no genomic intervention is needed to study a computer-mediated loop.

\section{Idea: A Closed-Loop Integration Benchmark}

\subsection{Central proposition}

The proposed CLOSED-LOOP-HUMAN-MACHINE-INTEGRATION-BENCHMARK (CHMI-Benchmark) evaluates integration as a vector of measurable properties. A useful coupling should improve task utility on held-out conditions, permit prediction and recovery when the system is perturbed, preserve the user's ability to veto, and allow safe disengagement. A high command rate alone is insufficient. The benchmark reports a Pareto profile over utility, safety, agency, independence, privacy, and security.

\subsection{Hypotheses}

H0 is the assistance-only null: an external device can improve immediate performance without improving closed-loop prediction, perturbation recovery, or machine-off transfer. H1 predicts that bidirectional feedback improves control error, latency, and recovery relative to matched unidirectional control. H2 predicts that joint human-machine adaptation improves acute performance but may increase dependence or sensitivity to decoder drift. H3 predicts that explicit veto and an independent stop preserve agency and safety. H4 predicts that held-out perturbations and machine-off transfer distinguish functional integration from calibration or assistance. H5 predicts a bandwidth-security tradeoff: richer signals create more opportunity for unintended inference and attack. H6 states a design boundary: a task-specific, reversible coupling is a safer near-term target than a species-level or genomic claim.

\subsection{What would count as evidence?}

The benchmark does not produce a scalar called ``hybridness.'' Evidence for a functional coupling requires a pre-registered gain that survives an unseen delay, drift, or task change; an improvement in predicting machine state from feedback; and no unacceptable loss in veto, safe shutdown, or independent operation. If feedback helps only during calibration, the appropriate conclusion is ``feedback-assisted training.'' If joint adaptation improves throughput but worsens machine-off transfer, the conclusion is ``performance-dependent coupling with governance risk.'' These interpretations are more informative than a yes/no species label.

\section{Formal Model and Metrics}

\subsection{A coupled dynamical system}

Let $x_h(t)$ denote a latent human motor or cognitive state, $x_m(t)$ the machine state, $y_h(t)$ the measured physiological signal, $u_m(t)$ the machine action, and $z_m(t)$ feedback returned to the user. A minimal observation model is
\begin{equation}
y_h(t)=g\bigl(x_h(t)\bigr)+\eta_h(t),
\label{eq:observation}
\end{equation}
where $g$ maps a latent state to a sensor signal and $\eta_h$ contains measurement noise, drift, missing channels, and physiological variation. The adaptive decoder and controller are represented by
\begin{equation}
u_m(t)=\pi_{\theta}\bigl(y_h(t),x_m(t)\bigr),
\label{eq:controller}
\end{equation}
where $\theta$ is updated according to the assigned adaptation regime. Feedback changes future human state according to
\begin{equation}
x_h(t+1)=f_h\bigl(x_h(t),u_m(t),z_m(t),w_h(t)\bigr),
\label{eq:humanstate}
\end{equation}
where $w_h(t)$ denotes task demand and disturbance. Equations \eqref{eq:observation}-\eqref{eq:humanstate} make the key point explicit: a decoder is not merely reading a fixed intention. Its output changes the task and therefore changes the future signal distribution.

\subsection{Utility and coupling}

For task $k$, define the observed performance vector
\begin{equation}
\mathbf{q}_k=\left[1-e_k,\; -\ell_k,\; r_k,\; v_k,\; d_k,\; -p_k,\; -a_k\right],
\label{eq:utility}
\end{equation}
where $e_k$ is normalized task error, $\ell_k$ is command latency, $r_k$ is perturbation recovery, $v_k$ is veto success, $d_k$ is machine-off transfer, $p_k$ is privacy leakage, and $a_k$ is attack success. Signs are chosen so larger values are preferable. The vector is retained rather than collapsed unless a decision maker specifies explicit weights.

The coupling component can be quantified with mutual information
\begin{equation}
I(X;U)=\sum_{x,u}p(x,u)\log\frac{p(x,u)}{p(x)p(u)},
\label{eq:mutualinformation}
\end{equation}
where $X$ is the intended state and $U$ is the machine action. Mutual information is not agency: a machine can predict a command while the user lacks a veto. It is therefore paired with attribution, veto, and independent-operation measures.

\subsection{Adaptation and drift}

The decoder drift between calibration time $t_0$ and evaluation time $t$ can be reported as
\begin{equation}
\Delta_{\theta}(t)=\left\|\theta(t)-\theta(t_0)\right\|_2,
\label{eq:drift}
\end{equation}
or, when parameter identities are not comparable, as the increase in held-out prediction error. The four adaptation regimes isolate whether improvement comes from user learning, machine learning, or their interaction. The design must log every model update and distinguish a planned update from an emergency safety intervention.

\subsection{Privacy and safety}

For a non-target attribute $S$, privacy leakage is operationalized as the best permitted inference accuracy from retained signals, or as the information quantity
\begin{equation}
L_{\mathrm{privacy}}=I(Y;S),
\label{eq:privacy}
\end{equation}
where $Y$ is the stored signal representation. The quantity is reported with the task utility because compression that preserves control may still preserve sensitive attributes. A safety run is acceptable only when every command satisfies hard bounds
\begin{equation}
\left\|u_m(t)\right\|\leq u_{\max},\qquad T_{\mathrm{stop}}\leq T_{\max},
\label{eq:safety}
\end{equation}
where $u_{\max}$ is a force, velocity, or workspace limit and $T_{\max}$ is the maximum permitted emergency-stop latency. These limits are enforced outside the adaptive controller.

\section{ExperimentDesign: Staged Protocol}

\subsection{Stage A: simulation-first}

The first stage uses synthetic physiological signals and openly licensed prerecorded traces where their licenses permit reuse. The simulator produces wearable inertial channels, electromyography-like signals, electroencephalography-like signals, and an implanted-compatible population-code emulator. No signal is presented as belonging to a real person. Signal drift, fatigue, dropout, delay, and task perturbation are generated from a fixed seed and a held-out parameter distribution.

Each cell of the factorial matrix is evaluated using at least 30 independent seeds. The generated users vary baseline skill, signal noise, adaptation rate, feedback sensitivity, and fatigue response. The evaluation set is held out by user parameters and by task perturbation, so a model cannot obtain a positive result by memorizing the training distribution. The simulation stage is also the first location for replay, spoofing, and unauthorized-inference tests.

\subsection{Stage B: hardware-in-the-loop}

The second stage connects the controller to a non-human fixture or robotic simulator. The fixture may include a virtual or physical robotic-arm surrogate, an exoskeleton actuator surrogate, and a tactile feedback transducer. No human contact or implant is needed. A deterministic safety supervisor limits workspace, force, velocity, command duration, and timeout. It has a physically or logically independent emergency stop and cannot be updated by the learning controller.

The fixture receives bounded sensor dropout, latency, calibration drift, and spoofed packets. A successful controller must degrade safely rather than silently producing unsafe actions. The test records the difference between a controller's confidence and its actual error. This is important because a coupled system that is wrong with high confidence can defeat both agency and safety even when its average accuracy appears acceptable.

\subsection{Stage C: low-risk human feasibility}

Only if Stages A and B satisfy the pre-registered safety thresholds should a separate ethics-reviewed feasibility study be considered. Its first tasks use surface electromyography, inertial sensing, eye tracking, or electroencephalography, together with a virtual or tabletop target. Assistance is low-force and reversible; feedback is non-nociceptive; participants can withdraw at any time. No implanted device, invasive stimulation, genomic intervention, or high-consequence decision is part of this protocol.

The human stage is not a shortcut to a clinical claim. It tests feasibility, learning, agency probes, data governance, and individual variability. Clinical populations and implanted systems require separate protocols with medical-device oversight, adverse-event monitoring, explant or disablement plans, and a new risk-benefit analysis.

\subsection{Factorial conditions}

Table \ref{tab:design} summarizes the conditions. The primary contrast holds signal quality, task demand, training time, and controller capacity constant while changing the feedback direction. A secondary contrast changes the adaptation regime. Governance is tested as a system factor rather than appended after performance measurement.

\begin{table}[t]
\caption{Core CHMI-Benchmark experimental factors}
\label{tab:design}
\centering
\scriptsize
\begin{tabular}{p{0.20\columnwidth}p{0.67\columnwidth}}
\toprule
Factor & Levels\\
\midrule
Interface & Wearable assistance; non-invasive physiological control; implanted-compatible emulator; bidirectional simulated feedback.\\
Feedback & None; delayed visual or task feedback; low-risk bidirectional sensory cue.\\
Adaptation & Frozen decoder; user-only; machine-only; joint.\\
Governance & Ordinary controller; transparent state display; explicit veto; local processing plus independent stop.\\
Stressors & Baseline; signal drift; channel dropout; feedback delay; task perturbation; sandbox spoofing.\\
\bottomrule
\end{tabular}
\end{table}

\subsection{Task suite}

The intention-decoding task estimates a target selection or continuous command from a physiological signal. Reach and grasp uses a virtual or fixture-mounted arm to acquire, transport, and release objects with changes in position and mass. A gait-assistance surrogate estimates gait phase and applies bounded assistance in simulation; it reports stability and an effort proxy, not a human metabolic result. A sensory-prediction task asks the user or model to predict device state from returned cues.

The agency task introduces rare controller deviations, with both announced probes and unannounced tests confined to low-consequence simulated actions. It measures whether the participant detects an unexpected action, vetoes it, and safely recovers. The machine-off task removes assistance after training, with a washout interval and a new but related task. Perturbation recovery changes target, delay, or available channels. The sandbox security task replays or spoofs inputs, and measures whether the supervisor blocks unsafe commands without accessing a production device.

\subsection{Component ablations and matched controls}

The benchmark includes controls that isolate the source of any apparent gain. A replay control presents the controller with a signal sequence generated by the same task but removes online coupling. A yoked-assistance control gives one condition the assistance trajectory produced for another condition, matched for average force or command count. A no-feedback control preserves the same actuation and training exposure while removing returned sensory information. These controls distinguish a benefit caused by extra practice or mechanical help from a benefit caused by a bidirectional loop.

The adaptation ablation is equally important. In the frozen-decoder condition, the machine parameters are fixed after calibration and only the user may learn. In the user-only condition, the decoder is fixed but the user receives feedback and can adapt their strategy. In the machine-only condition, the user's task policy is held as constant as the simulator permits while the decoder updates. In the joint condition, both policies update, with every update logged. The comparison must use the same update budget and the same number of training trials. Otherwise, joint adaptation could win simply because it receives more opportunities to fit noise.

\begin{table}[t]
\caption{Controls that separate assistance from integration}
\label{tab:controls}
\centering
\scriptsize
\begin{tabular}{p{0.22\columnwidth}p{0.60\columnwidth}}
\toprule
Control & Interpretation if the control matches the proposed system\\
\midrule
Replay & Improvement may depend on signal statistics or practice rather than online coupling.\\
Yoked assistance & Improvement may be mechanical assistance rather than adaptive human-machine coordination.\\
No feedback & Improvement that survives feedback removal is not evidence for a feedback mechanism.\\
Frozen decoder & Improvement measures user learning and exposes whether machine updates are necessary.\\
Machine-off transfer & Failure after removal indicates dependence, not necessarily unsafe dependence, but limits the integration claim.\\
Veto-only governance & A safety gain with unchanged utility supports agency-preserving control design.\\
\bottomrule
\end{tabular}
\end{table}

The proposed and control conditions should be paired at the level of a run whenever possible. Pairing reduces sensitivity to task difficulty and gives a direct counterfactual: what would the same intended trajectory have produced without feedback, without adaptation, or without online coupling? The analysis must retain both absolute and paired effects. A small average improvement with a large tail of failures is not acceptable for a physical interface.

\subsection{Calibration, drift, and feedback delay}

Each run begins with a calibration block that is not counted as evidence of integration. Calibration estimates the mapping from signal to task command and records signal quality, channel dropout, and initial latency. Evaluation blocks then contain familiar tasks, held-out target configurations, and a delayed-feedback block. The delayed block is not intended to mimic every biological delay; it is a controlled stressor that tests whether the user and decoder rely on a fragile timing coincidence.

Drift is injected in at least three ways: a slow rotation of the signal basis, a change in channel gain, and missing-channel bursts. The same drift schedule is applied across matched conditions. The controller may adapt only within its assigned regime. A model that changes parameters rapidly may appear robust in the short term while becoming difficult to audit. Therefore, the benchmark reports the number, magnitude, and timing of updates in addition to the final error.

Feedback is represented at three levels. A visual cue displays device state with a known delay. A task-level cue displays success or failure but not the intermediate machine state. A bidirectional cue maps device state to a sensory signal available for prediction. The cue map is fixed within a preregistered experiment, and its perceptual meaning is learned in a separate training block. This separation avoids interpreting a newly learned symbol as natural sensory embodiment. The correct scientific claim is that the feedback channel improves state estimation under the tested mapping.

\subsection{Agency probe implementation}

Agency probes are designed so that an unexpected action has low consequence and can be stopped. In the simulator, the controller occasionally applies a small, bounded deviation from the intended trajectory. In the fixture, the supervisor first verifies that the deviation remains inside the allowed workspace. The participant or user model receives a veto channel that is distinct from the ordinary command channel. Veto success is measured at the command level and at the physical-stop level; a late veto that still allows an unsafe movement is not counted as successful.

The benchmark also measures prospective prediction. Before a command is executed, the user or model selects the expected device state. Calibration is evaluated by the relationship between confidence and correct prediction. A system can be accurate without being predictable if its actions are opaque, and it can be predictable without being useful if it merely repeats a simple action. The combination of prediction, attribution, and veto gives a more concrete agency profile than a post-task questionnaire alone.

\subsection{Security and fault-injection matrix}

Security tests are conducted as controlled fault injection. Replay attacks resend an earlier valid signal. Spoofing changes a target channel while leaving other channels consistent. Dropout removes a subset of channels. Firmware or controller substitution is represented by a version mismatch in the fixture. Denial of service removes feedback or delays the command. Every injected event has a known start, end, and ground-truth label. The expected safe behavior is not always to stop: a controller may continue at a lower safe authority when feedback is missing, provided that this fallback is pre-specified and audited.

\begin{table}[t]
\caption{Fault-injection observations}
\label{tab:faults}
\centering
\scriptsize
\begin{tabular}{p{0.22\columnwidth}p{0.57\columnwidth}}
\toprule
Fault & Required observations\\
\midrule
Replay & Detection time, false acceptance, and whether old intent causes an unsafe action.\\
Spoofed channel & Detection time, command deviation, veto success, and supervisor response.\\
Channel dropout & Graceful degradation, confidence calibration, and fallback mode.\\
Version mismatch & Rejection or quarantine before actuation.\\
Feedback denial & Safe fallback, user awareness, and disengagement latency.\\
\bottomrule
\end{tabular}
\end{table}

These tests are not penetration testing of a deployed medical device. Their purpose is to expose the coupling between inference and actuation before human exposure. A security result is reported jointly with the task result: a controller that maintains throughput while accepting a replayed command has failed the benchmark's safety objective.

\section{Analysis, Reproducibility, and Expected Results}

\subsection{Estimands and statistical model}

For simulation and fixture stages, let $Y_{ijks}$ be an outcome for condition $i$, task $j$, generated user or participant $k$, and seed $s$. A hierarchical model is
\begin{equation}
Y_{ijks}=\beta_0+\beta_i+\gamma_j+b_k+c_s+\epsilon_{ijks},
\label{eq:hierarchical}
\end{equation}
where $\beta_i$ is the fixed effect of feedback, adaptation, governance, or stressor, $\gamma_j$ is a task effect, $b_k$ is a user effect, $c_s$ is a seed effect, and $\epsilon_{ijks}$ is residual variation. Interaction terms are pre-registered for feedback by adaptation and feedback by stressor. The principal contrast is bidirectional versus unidirectional feedback under matched conditions.

The analysis reports confidence intervals or Bayesian credible intervals, individual trajectories, worst-decile performance, calibration curves, missing-data patterns, and effect heterogeneity. Primary outcomes are task success, normalized error, command latency, perturbation recovery, veto success, emergency-stop latency, safe-disengagement latency, machine-off transfer, and decoder drift. Secondary outcomes are effort proxy, body-state prediction, workload, non-target inference, spoofing success, detection time, and subgroup variability. The multiplicity hierarchy is fixed before evaluation: utility and safety are primary; privacy, security, and independence are co-primary safeguards; exploratory metrics are labeled accordingly.

No scalar composite may hide a safety failure. A deployment recommendation requires a Pareto improvement or an explicitly justified tradeoff with no hard-threshold violation. The study should report a condition even if it performs poorly, because failure under drift or machine-off transfer is evidence about the boundary of integration.

\subsection{Decision rules}

A positive functional-coupling result requires five conditions. First, a preregistered improvement in held-out utility must be present. Second, the improvement must survive at least one unseen delay or drift condition. Third, veto success and independent emergency stopping must remain above their safety thresholds. Fourth, machine-off performance must not show unacceptable rebound errors or loss of basic capability. Fifth, the chosen data architecture must not create a material increase in non-target inference without an explicit consent and governance justification.

If feedback improves only calibration, the outcome is a training benefit. If joint adaptation raises throughput but worsens machine-off transfer or veto success, the outcome is performance-dependent coupling with governance risk. If a wearable, reversible system has the best safety-adjusted profile, the result favors that route for near-term deployment. None of these outcomes establishes a new species.

\subsection{Reproducible release}

The release package should contain the synthetic-data generator, fixed seed list, condition manifest, controller versions, model-update logs, metric definitions, safety logs, privacy threat model, security test scripts, and analysis code. The package should include a machine-readable record of exclusions and stopped runs. Raw physiological data, if collected in a later human stage, should be minimized, encrypted, separated from identity, processed locally where feasible, and deleted on request. Model updates must be versioned, reviewable, and disclosed to participants.

The result is reproducible only if the same signal generator, controller, seed, and threshold produce the same condition definitions. Randomness may be used inside a condition, but it must be recorded. A failed safety run must not be silently retried until it disappears from the summary. This requirement is particularly important when adaptive systems learn around a rare failure.

\section{Safety, Governance, and Boundaries}

\subsection{Safety as a control property}

The safety supervisor is a deterministic layer that limits action, speed, force, workspace, and timeout. The learning controller cannot modify these bounds. The stop channel is tested before every run and under sensor dropout. A system is stopped when it exceeds a bound, loses the safety channel, or produces an action whose confidence and context are incompatible. The protocol uses reversible operation and does not stimulate pain pathways.

This architecture addresses a basic asymmetry: a person may be harmed by one rare command even if average task accuracy is excellent. Thus the benchmark records maximum bounded action, supervisor intervention, stop latency, recovery time, and false-safe versus false-unsafe decisions. Safety failures are adverse events in the design record, not merely low-scoring observations.

\subsection{Agency, privacy, and cybersecurity}

Agency is operationalized through prediction, attribution, and veto. The benchmark measures intention-to-action latency, prediction of the machine's next action, correct attribution of self-generated versus controller-generated movement, successful veto, false veto, and recovery from a wrong command. User acceptance is not treated as evidence of agency. Transparency is tested by whether a state display improves calibration and veto rather than by whether users report liking it.

Physiological and neural signals are sensitive data. Data minimization, local processing, encryption, identity separation, role-based access, and deletion on request are part of the experiment. Equation \eqref{eq:privacy} turns the concern into a measurable quantity: the evaluator estimates the best permitted inference of a non-target attribute from stored representations. Secondary uses require new consent. Model updates require version records because a new decoder changes what the machine can infer.

Cybersecurity is biological safety when a networked controller can change physical action. Threats include replay, spoofed sensor input, malicious firmware, unauthorized inference, adversarial environmental input, and denial of service. Tests are performed in simulation and on the non-human fixture only. The protocol does not probe a clinical system or a production implant. NIST's Artificial Intelligence Risk Management Framework provides a useful governance vocabulary for validity, reliability, transparency, privacy, security, fairness, and accountability \cite{nistairmf}.

\subsection{Boundary of the species claim}

The protocol can test whether a human and computer jointly perform a specified task. It cannot test whether the pair constitutes a biological species, whether a machine is sentient, or whether human moral status has changed. It cannot infer that a cognitive augmentation is equivalent to a new mind. Genomic or heritable editing is explicitly excluded because it would introduce different medical, intergenerational, and governance risks. The word ``hybrid'' is therefore used only as shorthand for a functional human-machine coupling, with the qualification stated every time the claim is interpreted.

\section{Limitations and Discussion}

First, simulation and hardware-in-the-loop tests cannot reproduce every property of living neural tissue, including inflammation, electrode encapsulation, pain, fatigue, and changing clinical goals. An implanted-compatible emulator tests timing and signal statistics, not the safety or efficacy of implantation. Second, low-risk non-invasive validation cannot establish the feasibility of high-bandwidth intracortical feedback. Third, machine-off transfer is only one measure of independence; a user may choose dependence because the system is useful, and that choice is not itself evidence of harm or lack of agency.

Fourth, mutual information can be high for a control signal that is not consciously available, and agency can be preserved even when information is low if the user has a reliable veto. The benchmark therefore treats information, agency, and utility as distinct dimensions. Fifth, privacy leakage depends on which non-target attributes are evaluated. The threat model must be updated as new inference capabilities appear. Sixth, the proposed 30 seeds per simulation cell are a reproducibility minimum, not a universal power calculation. Human sample size must be derived from Stage B variance and reviewed by an ethics committee.

The broader implication is that integration is likely to be task-specific and layered. A person may use a powered exoskeleton, a speech decoder, or a sensory prosthesis without becoming a different biological kind of entity. The engineering question is whether the system is useful, robust, controllable, and governable. That question is both less sensational and more informative. It allows researchers to compare a reversible wearable with an implant, or a command decoder with a cognitive assistant, using common safeguards without pretending that they have the same scientific endpoint.

The benchmark also has a negative-result value. It may reveal that joint adaptation overfits the training environment, that feedback benefits vanish under delay, that privacy leakage rises faster than utility, or that machine-off transfer is consistently poor. Such findings would constrain claims about hybridization and direct future designs toward narrower, safer systems. A mature field should report these limits as prominently as its demonstrations.

\subsection{Interpretation matrix and deployment roadmap}

The same numerical improvement has different meanings at different integration levels. A wearable that lowers effort in a repeated gait task may be valuable even if it has no neural feedback channel. A non-invasive decoder that improves a communication task may be useful even if its signal does not become part of a body schema. An implanted interface may offer higher bandwidth while imposing surgical and maintenance risks that make its safety-adjusted utility lower for a particular application. The benchmark therefore maps outcomes to claims rather than ranking technologies by sophistication.

\begin{table}[t]
\caption{Claim boundaries for interpreting benchmark outcomes}
\label{tab:claims}
\centering
\scriptsize
\begin{tabular}{p{0.26\columnwidth}p{0.53\columnwidth}}
\toprule
Observed pattern & Defensible interpretation\\
\midrule
Higher task utility only & Assistance or decoding benefit; no evidence of closed-loop integration.\\
Higher utility and state prediction & Functional feedback coupling for the tested task and cue map.\\
Higher utility with lower veto or off-transfer & Coupling benefit accompanied by agency or dependence risk.\\
Robust gain under held-out drift & Evidence for adaptation that generalizes beyond calibration.\\
Low utility but strong privacy and safety & A governable design with insufficient current performance.\\
Any genomic change & A separate biological intervention, not an interface benchmark result.\\
\bottomrule
\end{tabular}
\end{table}

This matrix gives the Author a disciplined way to state a result. It prevents a high score on a cursor or reach task from being interpreted as a cognitive merger, and it prevents a lower-throughput but reversible device from being dismissed as scientifically uninteresting. It also makes the benchmark portable: the task can change from reach and grasp to speech, gait, or environmental sensing while the same claim boundaries remain in force.

The deployment roadmap has four gates. Gate one is computational: the controller must pass held-out simulation with known limits, reproducible seeds, calibration logs, and no hard-threshold safety violation. Gate two is fixture-based: the implementation must survive dropout, delay, drift, replay, and spoofing while the independent supervisor responds within the pre-specified time. Gate three is human feasibility: an ethics-reviewed, low-risk study verifies usability, agency, consent, and data deletion with reversible hardware. Gate four is application-specific translation: clinical, occupational, or consumer deployment receives its own benefit-risk analysis and does not inherit evidence from a different interface or population.

Each gate should produce a decision packet rather than only a mean score. The packet contains the outcome vector, worst-decile traces, safety incidents, model-update history, privacy evaluation, security faults, and a statement of what the data do not establish. A gate can be passed for one task and failed for another. In particular, passing a wearable gait task does not license a cognitive-augmentation claim, and passing a simulated implanted-compatible task does not license human implantation.

This staged logic also clarifies how to compare reversible and invasive systems. Reversibility has option value: a user can stop, replace, or update the device without permanent biological change. Invasive bandwidth can have value when a low-risk interface cannot achieve the needed function, but it must compensate for its added risk with a demonstrated gain in the same target task. The benchmark makes this comparison explicit by reporting safety-adjusted Pareto fronts rather than an unqualified technology ladder.

Finally, the roadmap treats users as adaptive partners rather than passive endpoints. A decoder update can change the meaning of a signal; a feedback mapping can change the user's strategy; and a shutdown policy can change trust and reliance. Versioning and disclosure are therefore experimental variables. If transparency improves veto calibration at a small throughput cost, the tradeoff may be favorable. If hidden adaptation improves accuracy while making actions less predictable, deployment should be delayed even when the average score is high. These are ordinary engineering decisions made unusually important by the fact that the actuator is coupled to a living person.

\subsection{Transfer, maintenance, and clinical translation}

An important test of integration is transfer across tasks and time. A controller trained on a reach task may learn a narrow mapping that fails when target geometry, object weight, or feedback delay changes. The benchmark therefore reserves a held-out task family rather than only held-out trials. A positive transfer result would mean that the coupled system has learned a reusable relation between intention, device state, and feedback. A negative result would identify a boundary of the claim and motivate task-specific controllers rather than a general-purpose hybrid assistant.

Long-term maintenance is part of the system model. Wearable sensors move relative to skin; electrodes and tissue interfaces can change; batteries and firmware require replacement; and the user's goals can change. A periodic recalibration burden should be reported as an operational cost, not hidden in a laboratory setup. The protocol should track time to restore safe performance, the amount of data needed for recalibration, and whether a model update changes the user's learned strategy. A system that is accurate only after frequent opaque retraining may be less governable than a lower-throughput system with predictable maintenance.

Translation also requires a clear separation between technical performance and clinical benefit. For a prosthesis, task success may be a meaningful intermediate endpoint, but comfort, skin health, infection risk, reliability, and quality of life determine whether the system helps. For a cognitive aid, a faster response may not be beneficial if it increases false confidence or exposes private information. The benchmark is therefore a preclinical comparison framework, not a substitute for clinical endpoints. Each application must select its own harm model and involve domain experts before deployment.

The most informative future experiment may be a longitudinal crossover in which participants use a reversible system across several tasks, with scheduled machine-off intervals and transparent model updates. Such a study would test whether the user learns a stable shared policy, whether the machine learns to predict the user's latent state, and whether governance features remain effective after familiarity grows. It should measure adaptation at multiple time scales and include planned withdrawal from the system. Withdrawal is not a punishment or a test of loyalty; it is a check that a person can retain meaningful choice about continued coupling.

This transfer perspective also changes what counts as a successful human-machine hybrid. A pair that performs one scripted action at high speed is a coupled actuator system. A pair that generalizes safely, exposes its uncertainty, supports veto, and remains useful after controlled changes is a stronger functional coupling. Even the stronger result remains a technological relation rather than proof of a new biological species. Keeping that distinction allows the research program to be ambitious about capability while precise about what its evidence means.

\section{Conclusion}

Current evidence supports useful human-machine couplings for bounded functions, including robotic-arm control through neural signals. It does not support the claim that humans and computers have formed a new biological species. This paper proposed CHMI-Benchmark to replace that broad claim with a causal and reproducible test: compare interface level, feedback, adaptation, and governance while measuring utility, perturbation recovery, agency, independence, privacy, security, and safe disengagement. The staged protocol begins with simulation, proceeds to a non-human hardware fixture, and reaches human validation only after ethics and safety review. Its central design principle is simple: a high-bandwidth interface is not a successful hybrid if the person cannot understand, veto, or safely leave it.

\balance
\begin{thebibliography}{00}
\bibitem{hochberg2012} L. R. Hochberg \textit{et al.}, ``Reach and grasp by people with tetraplegia using a neurally controlled robotic arm,'' \textit{Nature}, vol. 485, pp. 372--375, 2012, doi: 10.1038/nature11076.
\bibitem{hochberg2006} L. R. Hochberg \textit{et al.}, ``Neuronal ensemble control of prosthetic devices by a human with tetraplegia,'' \textit{Nature}, vol. 442, pp. 164--171, 2006, doi: 10.1038/nature04970.
\bibitem{gilja2011} V. Gilja \textit{et al.}, ``Challenges and opportunities for next-generation intracortically based neural prostheses,'' \textit{IEEE Transactions on Biomedical Engineering}, vol. 58, no. 7, pp. 1891--1899, 2011, doi: 10.1109/TBME.2011.2107553.
\bibitem{lebedev2006} M. A. Lebedev and M. A. L. Nicolelis, ``Brain-machine interfaces: past, present and future,'' \textit{Trends in Neurosciences}, vol. 29, no. 9, pp. 536--546, 2006, doi: 10.1016/j.tins.2006.07.004.
\bibitem{nistairmf} National Institute of Standards and Technology, \textit{Artificial Intelligence Risk Management Framework (AI RMF 1.0)}, NIST AI 100-1, 2023, doi: 10.6028/NIST.AI.100-1.
\end{thebibliography}

\end{document}
